% GENERATED FILE. Edit canonical lessons, then run npm run generate:books.
% canonical-source-hash: fnv1a64:a66b7655dde5db28
% canonical-lessons: UR-C414-gard, UR-C414-avara, UR-C414-bhaunkna, UR-C414-ghurrana, UR-C414-bhagana

\chapter{Guards and Street Dogs}
\label{ch:ur-guards-and-street-dogs}
\hlchaptermodality{414}

\begin{chapteropening}
\textbf{By the end of this chapter:} \emph{I can talk about the guard at the gate and the dogs in the lane.}

Complete the last lesson of chapter 414: I can talk about the guard at the gate and the dogs in the lane.
\end{chapteropening}

\section[\texorpdfstring{\ur{گارڈ} (gārḍ)}{گارڈ (gārḍ)}]{\ur{گارڈ} (gārḍ) — a guard}

\label{lesson:UR-C414-gard}



\begin{quote}
Before the new one: say the Urdu for a thatched roof, then the Urdu for to found (a town).
\end{quote}

\subsection*{You'll want to know: \ur{گارڈ}}
\textbf{\ur{گارڈ}} — \emph{gārḍ} — \textquotedblleft{}a guard\textquotedblright{}.

\begin{quote}
\ur{عمارت} \ur{کے} \ur{دروازے} \ur{پر} \ur{ایک} \ur{گارڈ} \ur{ہے۔}
\end{quote}

\begin{quote}
\emph{imārat ke darvāze par ek gārḍ hai.}
\end{quote}

\begin{quote}
There's a guard at the door of the building.
\end{quote}

Masculine or feminine, by the person meant.

\begin{cousinweb}
\textbf{\ur{گارڈ}} (\emph{gārḍ}) is borrowed from English \emph{guard}.
\end{cousinweb}

\begin{grammarlens}[title={What you've built}]
One more word for getting around.
\end{grammarlens}

\subsection*{Your turn}
\begin{itemize}
\raggedright
  \item \emph{Say it:} \emph{gārḍ}
  \item \emph{Say it:} \emph{gārḍ}, once more
  \item \emph{Say it:} \emph{imārat ke darvāze par ek gārḍ hai.} — There's a guard at the door of the building.
  \item \emph{Recall:} say the Urdu for a thatched roof, then the Urdu for to found (a town), then say \emph{gārḍ} again
\end{itemize}

\begin{tcolorbox}[breakable,colback=teal!4,colframe=teal!35!black,title={Before you move on}]
What does \ur{گارڈ} mean? (\textquotedblleft{}A guard\textquotedblright{}.) Say it once more.
\end{tcolorbox}
\section[\texorpdfstring{\ur{آوارہ} (āvārā)}{آوارہ (āvārā)}]{\ur{آوارہ} (āvārā) — stray; wandering}

\label{lesson:UR-C414-avara}



\begin{quote}
Before the new one: say the Urdu for to found (a town), then the Urdu for a guard.
\end{quote}

\subsection*{You'll want to know: \ur{آوارہ}}
\textbf{\ur{آوارہ}} — \emph{āvārā} — \textquotedblleft{}stray; wandering\textquotedblright{}.

\begin{quote}
\ur{گلی} \ur{میں} \ur{آوارہ} \ur{کتے} \ur{رات} \ur{کو} \ur{گھومتے} \ur{ہیں۔}
\end{quote}

\begin{quote}
\emph{galī meṅ āvārā kutte rāt ko ghūmte haiṅ.}
\end{quote}

\begin{quote}
Stray dogs roam the lane at night.
\end{quote}

\begin{cousinweb}
\textbf{\ur{آوارہ}} (\emph{āvārā}) is usually traced to Persian \emph{āvāra}, \textquotedblleft{}wandering, without a home\textquotedblright{}.
\end{cousinweb}

\begin{grammarlens}[title={What you've built}]
One more word for getting around.
\end{grammarlens}

\subsection*{Your turn}
\begin{itemize}
\raggedright
  \item \emph{Say it:} \emph{āvārā}
  \item \emph{Say it:} \emph{āvārā}, once more
  \item \emph{Say it:} \emph{galī meṅ āvārā kutte rāt ko ghūmte haiṅ.} — Stray dogs roam the lane at night.
  \item \emph{Recall:} say the Urdu for to found (a town), then the Urdu for a guard, then say \emph{āvārā} again
\end{itemize}

\begin{tcolorbox}[breakable,colback=teal!4,colframe=teal!35!black,title={Before you move on}]
What does \ur{آوارہ} mean? (\textquotedblleft{}Stray; wandering\textquotedblright{}.) Say it once more.
\end{tcolorbox}
\section[\texorpdfstring{\ur{بھونکنا} (bhauṅknā)}{بھونکنا (bhauṅknā)}]{\ur{بھونکنا} (bhauṅknā) — to bark}

\label{lesson:UR-C414-bhaunkna}



\begin{quote}
Before the new one: say the Urdu for a guard, then the Urdu for stray.
\end{quote}

\subsection*{You'll want to know: \ur{بھونکنا}}
\textbf{\ur{بھونکنا}} — \emph{bhauṅknā} — \textquotedblleft{}to bark\textquotedblright{}.

\begin{quote}
\ur{رات} \ur{کو} \ur{کتے} \ur{بھونکتے} \ur{ہیں۔}
\end{quote}

\begin{quote}
\emph{rāt ko kutte bhauṅkte haiṅ.}
\end{quote}

\begin{quote}
At night the dogs bark.
\end{quote}

Past (it …): \ur{وہ} \ur{کتا} \ur{دیر} \ur{تک} \ur{بھونکا}, \emph{voh kuttā der tak bhauṅkā}.

\begin{cousinweb}
\textbf{\ur{بھونکنا}} (\emph{bhauṅknā}) is inherited: CDIAL lists the Hindi cousins \emph{bhaũknā} and \emph{bhū̃knā}, \textquotedblleft{}to bark\textquotedblright{}, under Sanskrit \emph{búkkati}, \textquotedblleft{}barks\textquotedblright{} (CDIAL 9265).
\end{cousinweb}

\begin{grammarlens}[title={What you've built}]
One more word for getting around.
\end{grammarlens}

\subsection*{Your turn}
\begin{itemize}
\raggedright
  \item \emph{Say it:} \emph{bhauṅknā}
  \item \emph{Say it:} \emph{bhauṅknā}, once more
  \item \emph{Say it:} \emph{rāt ko kutte bhauṅkte haiṅ.} — At night the dogs bark.
  \item \emph{Recall:} say the Urdu for a guard, then the Urdu for stray, then say \emph{bhauṅknā} again
\end{itemize}

\begin{tcolorbox}[breakable,colback=teal!4,colframe=teal!35!black,title={Before you move on}]
What does \ur{بھونکنا} mean? (\textquotedblleft{}To bark\textquotedblright{}.) Say it once more.
\end{tcolorbox}
\section[\texorpdfstring{\ur{غرانا} (ghurrānā)}{غرانا (ghurrānā)}]{\ur{غرانا} (ghurrānā) — to growl}

\label{lesson:UR-C414-ghurrana}



\begin{quote}
Before the new one: say the Urdu for stray, then the Urdu for to bark.
\end{quote}

\subsection*{You'll want to know: \ur{غرانا}}
\textbf{\ur{غرانا}} — \emph{ghurrānā} — \textquotedblleft{}to growl\textquotedblright{}.

\begin{quote}
\ur{اجنبی} \ur{کو} \ur{دیکھ} \ur{کر} \ur{کتا} \ur{غرایا۔}
\end{quote}

\begin{quote}
\emph{ajnabī ko dekh kar kuttā ghurrāyā.}
\end{quote}

\begin{quote}
The dog growled when it saw the stranger.
\end{quote}

Past (it …): \ur{وہ} \ur{کتا} \ur{غرایا}, \emph{voh kuttā ghurrāyā}.

\begin{cousinweb}
\textbf{\ur{غرانا}} (\emph{ghurrānā}) is imitative: CDIAL lists the Hindi cousin \emph{gurrānā}, \textquotedblleft{}to growl\textquotedblright{}, under \emph{guragura}, a form it reconstructs from the sound (CDIAL 4207); Urdu writes it with \textbf{\ur{غ}}.
\end{cousinweb}

\begin{grammarlens}[title={What you've built}]
One more word for getting around.
\end{grammarlens}

\subsection*{Your turn}
\begin{itemize}
\raggedright
  \item \emph{Say it:} \emph{ghurrānā}
  \item \emph{Say it:} \emph{ghurrānā}, once more
  \item \emph{Say it:} \emph{ajnabī ko dekh kar kuttā ghurrāyā.} — The dog growled when it saw the stranger.
  \item \emph{Recall:} say the Urdu for stray, then the Urdu for to bark, then say \emph{ghurrānā} again
\end{itemize}

\begin{tcolorbox}[breakable,colback=teal!4,colframe=teal!35!black,title={Before you move on}]
What does \ur{غرانا} mean? (\textquotedblleft{}To growl\textquotedblright{}.) Say it once more.
\end{tcolorbox}
\section[\texorpdfstring{\ur{بھگانا} (bhagānā)}{بھگانا (bhagānā)}]{\ur{بھگانا} (bhagānā) — to chase off; to drive away}

\label{lesson:UR-C414-bhagana}



\begin{quote}
Before the new one: say the Urdu for to bark, then the Urdu for to growl.
\end{quote}

\subsection*{You'll want to know: \ur{بھگانا}}
\textbf{\ur{بھگانا}} — \emph{bhagānā} — \textquotedblleft{}to chase off; to drive away\textquotedblright{}.

\begin{quote}
\ur{گارڈ} \ur{نے} \ur{کتوں} \ur{کو} \ur{بھگایا۔}
\end{quote}

\begin{quote}
\emph{gārḍ ne kuttoṅ ko bhagāyā.}
\end{quote}

\begin{quote}
The guard chased the dogs away.
\end{quote}

Past, with \ur{نے} (he or she …): \ur{اس} \ur{نے} \ur{بلی} \ur{کو} \ur{بھگایا}, \emph{us ne billī ko bhagāyā}.

\begin{cousinweb}
\textbf{\ur{بھگانا}} (\emph{bhagānā}) is built on \textbf{\ur{بھاگنا}}, \textquotedblleft{}to run away\textquotedblright{}, which CDIAL lists under Sanskrit \emph{bhajyate}, \textquotedblleft{}is broken\textquotedblright{}, and its past form \emph{bhagna}, \textquotedblleft{}routed\textquotedblright{} (CDIAL 9361): this is the form for putting someone to flight.
\end{cousinweb}

\begin{grammarlens}[title={What you've built}]
One more word for getting around.
\end{grammarlens}

\subsection*{Your turn}
\begin{itemize}
\raggedright
  \item \emph{Say it:} \emph{bhagānā}
  \item \emph{Say it:} \emph{bhagānā}, once more
  \item \emph{Say it:} \emph{gārḍ ne kuttoṅ ko bhagāyā.} — The guard chased the dogs away.
  \item \emph{Recall:} say the Urdu for to bark, then the Urdu for to growl, then say \emph{bhagānā} again
\end{itemize}

\begin{tcolorbox}[breakable,colback=teal!4,colframe=teal!35!black,title={Before you move on}]
What does \ur{بھگانا} mean? (\textquotedblleft{}To chase off; to drive away\textquotedblright{}.) Say it once more.
\end{tcolorbox}
